\documentclass[a4paper,14pt]{extarticle}

% --- Шрифты и языки ---
\usepackage{fontspec} % поддержка любых шрифтов
\setmainfont{Times New Roman} % основной шрифт
\setmonofont{Courier New} % моноширинный шрифт с кириллицей
\newfontfamily\cyrillicfonttt{Courier New}
\usepackage{polyglossia}
\setmainlanguage{russian}
\setotherlanguage{english}

% --- Основные пакеты ---
\usepackage{graphicx}
\usepackage{caption}
\usepackage{setspace}
\usepackage{indentfirst}
\usepackage{titlesec}
\usepackage[left=30mm, right=15mm, top=20mm, bottom=20mm]{geometry}

% --- Основной текст ---
\setlength{\parindent}{1.25cm} % абзацный отступ
\onehalfspacing% межстрочный интервал

% --- Подписи к рисункам ---
\captionsetup{
	labelsep=period,
	justification=centering,
	font=normalsize
}

% --- Заголовки разделов (1 уровень) ---
\titleformat{\section}
{\bfseries\centering\fontsize{16}{19}\selectfont}
{\thesection}{1em}{}
\titlespacing*{\section}{0pt}{2\baselineskip}{1.5\baselineskip}

% --- Заголовки подразделов (2 уровень) ---
\titleformat{\subsection}
{\bfseries\fontsize{15}{18}\selectfont}
{\thesubsection}{1em}{}
\titlespacing*{\subsection}{0pt}{1.5\baselineskip}{1\baselineskip}

\sloppy


\begin{document}

\begin{titlepage}
	\centering
	{\large Донской Государственный Технический Университет}\\[4cm]

	{\bfseries\LARGE Лабораторные работы\\[0.5cm]
	{\large по дисциплине <<Разработка пользовательского инферфейса>>}\\[6cm]

	\begin{flushright}
		Выполнил: студент группы ВИАС21\\
		Илья Бондаренко
	\end{flushright}

	\vfill

	Ростов-на-Дону --- 2025
\end{titlepage}

\tableofcontents
\newpage

\section{Лабораторная работа №1}

\subsection{Создание структуры проекта}
По примеру из задания, был создан простейший HTML-документ.

\begin{figure}[h]
	\centering
	\includegraphics[width=1\textwidth]{./images/1/1.1_code.jpg}
	\caption{Код базовой структуры HTML-документа}
\end{figure}

\subsection{Форматирование. Работа с текстом}
Целью данного задания было испытание возможностей оформления веб-страницы с помощью тэгов HTML их аттрибутов. В конце концов получился документ по химической тематике, каждая строка которого была оформлена по разному.

\begin{figure}[ht]
	\centering
	\includegraphics[width=1\textwidth]{./images/1/1.2_code.jpg}
	\caption{Фрагмент кода для веб-страницы с разнообразным оформлением}
\end{figure}

\begin{figure}[ht]
	\centering
	\includegraphics[width=1\textwidth]{./images/1/1.2_result.jpg}
	\caption{Веб страница с разнообразным оформлением}
\end{figure}

\subsection{Форматирование стихотворения}
В этом задании потребовалось выбрать произвольное стихотворение, вывести его на веб-страницу, а также основную информацию об его авторе. После чего, пользуясь возможностями HTML выделить ключевые слова стихотворения.

\begin{figure}[ht]
	\centering
	\includegraphics[width=1\textwidth]{./images/1/1.3_code.png}
	\caption{Фрагмент кода для веб-страницы со стихотворением}
\end{figure}

\begin{figure}[ht]
	\centering
	\includegraphics[width=1\textwidth]{./images/1/1.3_result.jpg}
	\caption{Фрагмент веб-страницы со стихотворением}
\end{figure}

\newpage

\subsection{Форматирование текста по образцу}
В последнем задании этой лабораторной работы потребовалось оформить
полученный текст семантически правильно, соблюдая структуру оформления.

\begin{figure}[ht]
	\centering
	\includegraphics[width=1\textwidth]{./images/1/1.4_code.jpg}
	\caption{Фрагмент кода для веб-страницы об Алтае}
\end{figure}

\begin{figure}[ht]
	\centering
	\includegraphics[width=1\textwidth]{./images/1/1.4_result.jpg}
	\caption{Фрагмент веб-страницы об Алтае}
\end{figure}

\clearpage
\section{Лабораторная работа №2}

\subsection{Объявление стилей}
Рассмотрены способы подключения CSS: внутри тегов, в блоке <style> и во внешнем файле.
Изучены классы и идентификаторы.

\subsection{Списки и ссылки}
Созданы вложенные списки с разными типами нумерации.
Добавлены ссылки и применены hover-эффекты.

\begin{figure}[ht]
	\centering
	\includegraphics[width=1\textwidth]{./images/1/2.1_result.jpg}
	\caption{Фрагмент веб-страницы}
\end{figure}

\begin{figure}[h]
	\centering
	\includegraphics[width=1\textwidth]{./images/1/2.1_html.jpg}
	\caption{Фрагмент кода для веб-страницы}
\end{figure}

\clearpage

\subsection{Блоки и изображения}
Изучена работа с контейнером <div> и свойством position.
В блоки вставлены изображения.

\begin{figure}[ht]
	\centering
	\includegraphics[width=1\textwidth]{./images/1/2.3_result.jpg}
	\caption{Фрагмент веб-страницы}
\end{figure}

\begin{figure}[h]
	\centering
	\includegraphics[width=1\textwidth]{./images/1/2.3_html.jpg}
	\caption{Фрагмент кода для веб-страницы}
\end{figure}

\begin{figure}[h]
	\centering
	\includegraphics[width=1\textwidth]{./images/1/2.2_css.jpg}
	\caption{Код CSS-документа}
\end{figure}

\clearpage

\subsection{Навигация}
Созданы страницы index.html и contacts.html.
Реализована перелинковка, использованы протоколы tel:, mailto:, а также ссылка на PDF с атрибутом download.

\begin{figure}[ht]
	\centering
	\includegraphics[width=1\textwidth]{./images/1/2.4_result.jpg}
	\caption{Фрагмент веб-страницы}
\end{figure}

\begin{figure}[h]
	\centering
	\includegraphics[width=1\textwidth]{./images/1/2.4_html.jpg}
	\caption{Фрагмент кода для веб-страницы}
\end{figure}

\clearpage

\subsection{Верстка статьи}
Сверстана статья <<Осенние путешествия по России>>.
Список городов оформлен как гиперссылки, создана отдельная страница с достопримечательностями.

\begin{figure}[ht]
	\centering
	\includegraphics[width=1\textwidth]{./images/1/2.5_result.jpg}
	\caption{Фрагмент веб-страницы}
\end{figure}

\begin{figure}[ht]
	\centering
	\includegraphics[width=1\textwidth]{./images/1/2.5_result_0.jpg}
	\caption{Фрагмент вложенной веб-страницы}
\end{figure}

\begin{figure}[h]
	\centering
	\includegraphics[width=1\textwidth]{./images/1/2.5_html.jpg}
	\caption{Фрагмент кода для веб-страницы}
\end{figure}

\clearpage

\section{Лабораторная работа №3}

\subsection{Создание простейшей таблицы}
Создан файл HTML и подключён CSS. Оформлена базовая таблица с тегами <table>, <tr>, <td>. Добавлены стили для рамок и отступов, проверено отображение в браузере.

\begin{figure}[ht]
	\centering
	\includegraphics[width=1\textwidth]{./images/1/3.1_result.jpg}
	\caption{Фрагмент веб-страницы}
\end{figure}

\begin{figure}[h]
	\centering
	\includegraphics[width=0.5\textwidth]{./images/1/3.1_code.jpg}
	\caption{Фрагмент кода для веб-страницы}
\end{figure}

\subsection{Работа с рамками и отступами}
Настроены рамки ячеек: оставлены только правая и левая. Внутренние отступы ячеек заданы с помощью padding: сверху-снизу 5px, справа-слева 10px.

\begin{figure}[ht]
	\centering
	\includegraphics[width=1\textwidth]{./images/1/3.2_result.jpg}
	\caption{Фрагмент веб-страницы}
\end{figure}

\begin{figure}[h]
	\centering
	\includegraphics[width=0.5\textwidth]{./images/1/3.2_code.jpg}
	\caption{Фрагмент кода для веб-страницы}
\end{figure}

\clearpage

\subsection{Ячейки-заголовки}
Теги <td> в заголовках заменены на <th>. Для <th> установлены рамка и внутренние отступы.

\begin{figure}[ht]
	\centering
	\includegraphics[width=1\textwidth]{./images/1/3.3_result.jpg}
	\caption{Фрагмент веб-страницы}
\end{figure}

\begin{figure}[h]
	\centering
	\includegraphics[width=0.5\textwidth]{./images/1/3.3_code.jpg}
	\caption{Фрагмент кода для веб-страницы}
\end{figure}

\subsection{Объединение ячеек}
Использованы атрибуты colspan и rowspan для объединения ячеек по горизонтали и вертикали. Проверено отображение объединённых ячеек в таблице.

\begin{figure}[ht]
	\centering
	\includegraphics[width=0.5\textwidth]{./images/1/3.4_result.jpg}
	\caption{Фрагмент веб-страницы}
\end{figure}

\begin{figure}[h]
	\centering
	\includegraphics[width=0.5\textwidth]{./images/1/3.4_html.jpg}
	\caption{Фрагмент HTML кода для веб-страницы}
\end{figure}

\begin{figure}[h]
	\centering
	\includegraphics[width=0.5\textwidth]{./images/1/3.4_css.jpg}
	\caption{Фрагмент CSS кода для веб-страницы}
\end{figure}

\clearpage

\subsection{Создание таблицы с объединением}
Создана таблица с объединёнными заголовочными ячейками. Проверена корректная структура и отображение в браузере.

\begin{figure}[ht]
	\centering
	\includegraphics[width=0.5\textwidth]{./images/1/3.5_result.jpg}
	\caption{Фрагмент веб-страницы}
\end{figure}

\begin{figure}[h]
	\centering
	\includegraphics[width=0.5\textwidth]{./images/1/3.5_code.jpg}
	\caption{Фрагмент кода для веб-страницы}
\end{figure}

\subsection{Таблица с внутренними отступами кратными 5}
Настроены внутренние отступы для всех ячеек согласно заданию, кратные 5px.

\begin{figure}[ht]
	\centering
	\includegraphics[width=1\textwidth]{./images/1/3.7_result.jpg}
	\caption{Фрагмент веб-страницы}
\end{figure}

\begin{figure}[h]
	\centering
	\includegraphics[width=0.5\textwidth]{./images/1/3.7_code.jpg}
	\caption{Фрагмент кода для веб-страницы}
\end{figure}

\clearpage

\subsection{Полная таблица с thead, tbody и tfoot}
Создан HTML-документ с таблицей 5$\times$5. Верхняя строка обёрнута в thead с тегами <th>, средние строки — tbody, нижняя строка — tfoot с объединённой ячейкой и текстом «Всего: 3». CSS задаёт границы, отступы, ширину таблицы, цвет фона для thead и tfoot, текст в tfoot прижат к правому краю.

\begin{figure}[ht]
	\centering
	\includegraphics[width=1\textwidth]{./images/1/3.6_result.jpg}
	\caption{Фрагмент веб-страницы}
\end{figure}

\begin{figure}[h]
	\centering
	\includegraphics[width=0.5\textwidth]{./images/1/3.6_code.jpg}
	\caption{Фрагмент кода для веб-страницы}
\end{figure}

\clearpage

\subsection{Таблица с эффектами hover}
Применён селектор :nth-child () для горизонтальных полос «зебра». Настроены вертикальные полосы и горизонтальные разделители через border-bottom. Добавлен эффект выделения строки при наведении курсора с помощью tr:hover.

\begin{figure}[ht]
	\centering
	\includegraphics[width=1\textwidth]{./images/1/3.10_result.jpg}
	\caption{Фрагмент веб-страницы}
\end{figure}

\begin{figure}[h]
	\centering
	\includegraphics[width=0.5\textwidth]{./images/1/3.10_code.jpg}
	\caption{Фрагмент кода для веб-страницы}
\end{figure}

\clearpage

\section{Лабораторная работа №4}

\subsection{Задание №1}

\begin{figure}[ht]
	\centering
	\includegraphics[width=1\textwidth]{./images/4/1_result.png}
	\caption{Фрагмент веб-страницы}
\end{figure}

\begin{figure}[h]
	\centering
	\includegraphics[width=1\textwidth]{./images/4/1_html_code.png}
	\caption{Фрагмент HTML кода для веб-страницы}
\end{figure}

\begin{figure}[h]
	\centering
	\includegraphics[width=1\textwidth]{./images/4/1_css_code.png}
	\caption{Фрагмент CSS кода для веб-страницы}
\end{figure}

\clearpage

\subsection{Задание №2}

\begin{figure}[ht]
	\centering
	\includegraphics[width=1\textwidth]{./images/4/2_result.png}
	\caption{Фрагмент веб-страницы}
\end{figure}

\begin{figure}[h]
	\centering
	\includegraphics[width=1\textwidth]{./images/4/2_html_code.png}
	\caption{Фрагмент HTML кода для веб-страницы}
\end{figure}

\begin{figure}[h]
	\centering
	\includegraphics[width=1\textwidth]{./images/4/2_css_code.png}
	\caption{Фрагмент CSS кода для веб-страницы}
\end{figure}

\clearpage

\subsection{Задание №3}

\begin{figure}[ht]
	\centering
	\includegraphics[width=0.5\textwidth]{./images/4/3_result.png}
	\caption{Фрагмент веб-страницы}
\end{figure}

\begin{figure}[h]
	\centering
	\includegraphics[width=1\textwidth]{./images/4/3_html_code.png}
	\caption{Фрагмент HTML кода для веб-страницы}
\end{figure}

\begin{figure}[h]
	\centering
	\includegraphics[width=1\textwidth]{./images/4/3_css_code.png}
	\caption{Фрагмент CSS кода для веб-страницы}
\end{figure}

\clearpage

\subsection{Задание №4}

\begin{figure}[ht]
	\centering
	\includegraphics[width=1\textwidth]{./images/4/4_result.png}
	\caption{Фрагмент веб-страницы}
\end{figure}

\begin{figure}[h]
	\centering
	\includegraphics[width=1\textwidth]{./images/4/4_html_code.png}
	\caption{Фрагмент HTML кода для веб-страницы}
\end{figure}

\begin{figure}[h]
	\centering
	\includegraphics[width=1\textwidth]{./images/4/4_css_code.png}
	\caption{Фрагмент CSS кода для веб-страницы}
\end{figure}

\clearpage

\subsection{Задание №5}

\begin{figure}[ht]
	\centering
	\includegraphics[width=1\textwidth]{./images/4/5_result.png}
	\caption{Фрагмент веб-страницы}
\end{figure}

\begin{figure}[h]
	\centering
	\includegraphics[width=1\textwidth]{./images/4/5_html_code.png}
	\caption{Фрагмент HTML кода для веб-страницы}
\end{figure}

\begin{figure}[h]
	\centering
	\includegraphics[width=1\textwidth]{./images/4/5_css_code.png}
	\caption{Фрагмент CSS кода для веб-страницы}
\end{figure}

\clearpage

\section {Лабораторная работа №5}
\subsection {Задание}

\begin{figure}[ht]
	\centering
	\includegraphics[width=1\textwidth]{./images/5/1_result.png}
	\caption{Фрагмент веб-страницы №1}
\end{figure}

\begin{figure}[ht]
	\centering
	\includegraphics[width=1\textwidth]{./images/5/1.1_result.png}
	\caption{Фрагмент веб-страницы №2}
\end{figure}

\begin{figure}[h]
	\centering
	\includegraphics[width=1\textwidth]{./images/5/html.png}
	\caption{Фрагмент HTML кода для веб-страницы}
\end{figure}

\begin{figure}[h]
	\centering
	\includegraphics[width=1\textwidth]{./images/5/css.png}
	\caption{Фрагмент CSS кода для веб-страницы}
\end{figure}
\clearpage

\section {Лабораторная работа №6}
\subsection {Задание 1}

\begin{figure}[ht]
	\centering
	\includegraphics[width=0.8\textwidth]{./images/6/1_result.png}
	\caption{Фрагмент веб-страницы}
\end{figure}

\begin{figure}[h]
	\centering
	\includegraphics[width=1\textwidth]{./images/6/1_html.png}
	\caption{Фрагмент HTML кода для веб-страницы}
\end{figure}

\begin{figure}[h]
	\centering
	\includegraphics[width=1\textwidth]{./images/6/1_css.png}
	\caption{Фрагмент CSS кода для веб-страницы}
\end{figure}
\clearpage

\subsection {Задание 2}

\begin{figure}[ht]
	\centering
	\includegraphics[width=0.8\textwidth]{./images/6/2_result.png}
	\caption{Фрагмент веб-страницы}
\end{figure}

\begin{figure}[h]
	\centering
	\includegraphics[width=0.8\textwidth]{./images/6/2_html.png}
	\caption{Фрагмент HTML кода для веб-страницы}
\end{figure}
\clearpage

\end{document}